\chapter{Structural Functionalism: M. N. Srinivas}

\section{M. N. Srinivas: Background \& Mentors}

M. N. Srinivas (full name \textbf{Mysore Narasimhachar Srinivas}) is the first person to introduce structural functionalism to study India (under Radcliffe-Brown). He is called the \textbf{Doyen of Indian sociology} (Satish Deshpande's term), having contributed to family, politics, caste, village and religion. Other structural functionalists: \textbf{S. C. Dube} (critiqued Srinivas from within the framework, but Marxist-inclined), \textbf{I. P. Desai} (family/marriage/kinship), and \textbf{McKim Marriott} (some place him in the civilizational perspective).

\begin{KeyConcept}
Srinivas wrote an article titled \textbf{``Professor G. S. Ghurye and I''} --- a relationship of rivalry. Ghurye was his PhD supervisor, yet Srinivas never wanted to be a field worker; he wanted to be a theoretician like Max Weber. Ghurye made him a field worker.
\end{KeyConcept}

\begin{DataFact}
Srinivas did \textbf{two PhDs}: first under G. S. Ghurye (Bombay University); second, a \textbf{D.Phil at Oxford} under \textbf{A. R. Radcliffe-Brown} (and, after Radcliffe-Brown died, \textbf{Evans-Pritchard}). His three gurus are Ghurye, Radcliffe-Brown and Evans-Pritchard.
\end{DataFact}

\begin{ComparisonBox}
\begin{itemize}
\item \textbf{British tradition} of structural functionalism (Radcliffe-Brown, Evans-Pritchard) is \textbf{field-based}.
\item \textbf{French} (Durkheim) and \textbf{American} (Talcott Parsons) traditions are \textbf{theory-based} --- Durkheim and Parsons only theorised; they never did field study.
\end{itemize}
\end{ComparisonBox}

\begin{QuickRevision}
\begin{itemize}
\item Srinivas = first to introduce structural functionalism to India; ``Doyen of Indian sociology''.
\item Other SF scholars: S. C. Dube, I. P. Desai, McKim Marriott.
\item Two PhDs: Ghurye, then Radcliffe-Brown / Evans-Pritchard (Oxford D.Phil).
\item British SF = field-based; French/American = theory-based.
\end{itemize}
\end{QuickRevision}

\section{Why Coorg (Kurg)?}

Coorg (Kurg) is called the \textbf{``Scotland of India''} --- an isolated hilly region in the south-western part of Karnataka. Srinivas studied the \textbf{Coorgs} (the inhabitants) for two reasons:
\begin{itemize}
\item Ghurye wanted to check his diffusion theory: if even this isolated pocket shows Hindu character, then Vedic Hindu Aryan culture has diffused all over India.
\item Ghurye believed the Coorgs and Egyptian culture share features --- both preserve dead bodies (tombs).
\end{itemize}

\begin{ExampleBox}
Srinivas being from Mysore (Karnataka) also explains the location --- the geographical location of one's field is influenced by one's social location.
\end{ExampleBox}

\begin{QuickRevision}
\begin{itemize}
\item Coorg = ``Scotland of India'', isolated Karnataka region.
\item Two reasons: check Ghurye's diffusion theory; check Egyptian (tomb) connection.
\end{itemize}
\end{QuickRevision}

\section{The Coorg Study (1952)}

This study was the \textbf{beginning of the field-work tradition in India} --- a paradigm shift from \textbf{book view to field view}. Srinivas first relied on British records (Coorg was invaded by Lingayat Rajas in the 9th century AD; the British attacked the Lingayats in 1834 and conducted a census in 1881).

\begin{DataFact}
The census found the Coorgs \textbf{claimed Kshatriya status}, but the British gave them \textbf{Shudra status} (influenced by the racial theory of caste --- dark isolated tribals could not be Kshatriya). The Coorgs demanded Kshatriya status; in \textbf{1926} the Bombay High Court ruled they were \textbf{not Shudras}.
\end{DataFact}

The Coorgs had been the military/aristocracy under the Lingayat (Shiva-worshipping) dynasty, hence they claimed Kshatriya status. In the 1941 census the Coorgs had subdivisions --- \textbf{Amma Coorg} (Brahmin Coorg, about 600) and \textbf{Sanna Coorg} (Kshatriya Coorg). The \textbf{Kaveri Purana} (a local Purana) was written to justify their Kshatriya status --- tracing descent from the marriage of a disciple of sage Agastya with a Brahmin girl, after which the Kaveri river began flowing.

\begin{KeyConcept}
Srinivas fused structural functionalism with \textbf{historical sociology} --- unlike Western structural functionalism, which is \textbf{ahistorical}.
\end{KeyConcept}

\begin{QuickRevision}
\begin{itemize}
\item Coorg study = beginning of field-work tradition (paradigm shift from book to field view).
\item Coorgs claimed Kshatriya, British gave Shudra; Bombay HC (1926) ruled not Shudras.
\item Kaveri Purana justified Kshatriya status; Srinivas fused SF with historical sociology.
\end{itemize}
\end{QuickRevision}

\section{The Four Levels of Hinduism}

Hinduism has \textbf{no founder} (unlike Buddhism, Sikhism, Islam, Christianity). Srinivas believed Hinduism is \textbf{amorphous and complex}, a set of very diverse beliefs that varies from region to region. He spoke of \textbf{four levels of Hinduism}:
\begin{itemize}
\item \textbf{All-India Sanskritic Hinduism} --- Brahmins, Vedas, rituals, Upanayana (what we normally think of as Hinduism).
\item \textbf{Local Hinduism} --- local deities not found in all-India Hinduism (e.g.\ Murugan, Karni Mata).
\end{itemize}

\begin{KeyConcept}
``Edge effect'' (a concept Srinivas used, attributed to Radcliffe-Brown) and ``spread'': a smaller area has more homogeneous customs; as the area grows, heterogeneous elements appear.
\end{KeyConcept}

When a region gets ``Hinduised'', it means all-India Sanskritic Hinduism is absorbing local Hinduism, or local Hinduism is borrowing from it (a two-way process).

\begin{QuickRevision}
\begin{itemize}
\item Hinduism has no founder; it is amorphous and complex, regionally varied.
\item Four levels: all-India Sanskritic vs local (Murugan, Karni Mata); edge effect / spread.
\end{itemize}
\end{QuickRevision}

\section{The Village Structure of Coorg}

Srinivas found five important things in the Coorg village structure:
\begin{itemize}
\item \textbf{Oka} --- the joint patrilineal family system (property passed from father to son).
\item \textbf{Ritual purity and pollution} --- ``madi'' (pure) and ``pole'' (impure); birth, death, bodily emissions are impure; the cobra is a sacred shrine for ancestor worship.
\item \textbf{Caste structure} --- Brahmin Coorg, Kshatriya Coorg, and local castes (Banna, Meda, etc.).
\item \textbf{Interdependence between castes} --- dependence on upper castes (Brahmins) in ritual contexts (weddings), and on low castes in non-ritual contexts (e.g.\ a plumber to fix a toilet).
\item \textbf{Festivals} --- a time of social solidarity, bringing all castes together (collective effervescence).
\end{itemize}

\begin{QuickRevision}
\begin{itemize}
\item Coorg village structure: Oka (patrilineal joint family), purity/pollution (madi/pole), caste structure, interdependence, festivals.
\end{itemize}
\end{QuickRevision}

\section{Structural Functionalism Applied}

\begin{KeyConcept}
Structural functionalism looks at the units of social structure and how each contributes to the \textbf{maintenance of the system}. Radcliffe-Brown: ``a social structure is any structure that performs certain functions'' (system maintenance). Man is a biological being until initiated into rituals, becoming a \textbf{social personality}.
\end{KeyConcept}

Seen through this lens, the Coorg features maintain the system: \textbf{Oka} (jointness) maintains the units; \textbf{purity/pollution} produces value consensus (Parsons); \textbf{caste structure} creates interdependence; \textbf{festivals} bind people into an integrated whole.

\begin{ExampleBox}
Robert Merton's \textbf{rain dance} (among an African community): rain dance does not bring rain, but it brings solidarity and integrates the people --- its \textbf{latent function} (unanticipated consequence).
\end{ExampleBox}

\begin{QuickRevision}
\begin{itemize}
\item SF: each unit performs a function for system maintenance (Radcliffe-Brown).
\item Coorg units: Oka, purity/pollution, caste, festivals all maintain the system.
\item Merton's rain dance = latent function (solidarity, not rain).
\end{itemize}
\end{QuickRevision}

\section{Sanskritization}

\begin{KeyConcept}
\textbf{Sanskritization} (coined in \emph{Religion and Society Among the Coorgs of South India}): a process characterised by the \textbf{cultural imitation of the dvija (twice-born --- B, K, V)} by a low caste, tribe, or other group. E.g.\ adopting vegetarianism, the sacred thread, learning Sanskrit, teetotalism.
\end{KeyConcept}

It is essentially the application of \textbf{Merton's reference group theory}: the reference group of the low caste/tribal was the upper caste (dvija). Till the late 1940s people looked up to Brahmins; today the reference groups are secular (scientists, IAS officers).

\begin{ExampleBox}
The Coorgs Sanskritized themselves (worshipping Shiva, claiming Kshatriya status), producing Brahmin Coorg and Kshatriya Coorg within. Ancient examples: kings claimed Kshatriya status after invading; Maharishi Valmiki (Shudra became Maharishi); Buddha (Kshatriya claimed Brahminical status). Kanishka (``devaputra'') and Ashoka (``Devanampiya Piyadassi'') claimed status.
\end{ExampleBox}

\begin{CriticalPoint}
Sanskritization brings \textbf{social cohesion} (low castes become like high castes), but the high caste will not acknowledge them as equal. It changes \textbf{positions} but not the \textbf{system} --- the caste system is maintained despite mobility (moving equilibrium).
\end{CriticalPoint}

When tribes get Hinduised, egalitarian people absorb inequality --- but even this inequality is seen as functional (different groups dependent on each other).

\begin{QuickRevision}
\begin{itemize}
\item Sanskritization = cultural imitation of the dvija by low caste/tribe/group (Merton's reference group theory).
\item Brings social cohesion but not equality; changes positions, not the system.
\item Tribes getting Hinduised = egalitarians absorbing inequality.
\end{itemize}
\end{QuickRevision}

\section{The Remembered Village}

Srinivas's second field study was in \textbf{Rampura} (a pseudonym for \textbf{Kodagahalli}, Karnataka), conducted in 1948, days after Gandhi's assassination, on the advice of Evans-Pritchard. He used a pseudonym (like Shakespeare writing ``Hamlet'' for Helsing\o r) to protect villagers' privacy.

\begin{DataFact}
The social composition of Rampura (1948): \textbf{19 caste groups + 1 Muslim family}. The dominant caste was the \textbf{Vokkaliga (Okaliga)} --- a peasant, land-owning group with maximum numbers. Srinivas stayed 11 months and kept field notes.
\end{DataFact}

\begin{ExampleBox}
In the 1970s his field notes \textbf{were burnt} (during Vietnam-war protests in America, students set fire to STD booths). \textbf{Sol Tax} (an American anthropologist) advised him to recollect and rewrite from memory --- hence the book's name, \emph{The Remembered Village}.
\end{ExampleBox}

Concepts coined in this book: \textbf{dominant caste}, \textbf{vote bank}, and (reused) \textbf{Sanskritization}.

\begin{itemize}
\item Observations: Brahmins and Lingayats were proponents of law and justice; Vokkaligas resolved disputes; a Muslim group performed ``magic''/healing (hakeem); the \textbf{Jajmani system} established interdependence (low castes gave goods/services to upper castes for financial support).
\end{itemize}

\begin{QuickRevision}
\begin{itemize}
\item Rampura = pseudonym for Kodagahalli; 19 castes + 1 Muslim family; Vokkaliga dominant.
\item Field notes burnt (1970s, Vietnam protests); Sol Tax advised rewriting from memory.
\item Coined: dominant caste, vote bank; observations incl.\ Jajmani system.
\end{itemize}
\end{QuickRevision}

\section{Dominant Caste \& the Two Hierarchies}

\begin{KeyConcept}
\textbf{Dominant caste} = a caste with preponderance of \textbf{numerical strength, economic wealth and political power} (three capitals: social, economic, cultural/symbolic --- easily transmissible and interchangeable).
\end{KeyConcept}

\begin{ComparisonBox}
\begin{itemize}
\item \textbf{Ritual hierarchy} (book view, all-India): Brahmin supreme, Shudra inferior, on the basis of \textbf{purity and pollution}.
\item \textbf{Secular hierarchy} (field view, local): based on \textbf{land, wealth, power and numbers}.
\end{itemize}
\end{ComparisonBox}

In Rampura the Vokkaligas were dominant in the \textbf{secular} hierarchy; Brahmins were superior in the \textbf{ritual} hierarchy --- status is \textbf{contextual}. Even a dominant caste wants Kshatriya status, so it claims ritual superiority too.

\begin{ExampleBox}
The \textbf{Yadavs} of U.P. became a dominant caste (numbers, land, ministers, vote bank) --- but they first had to Sanskritize themselves, creating myths (descent from Yadu) to claim Kshatriya status.
\end{ExampleBox}

\begin{PYQLink}
The teacher cites the PYQ ``Salient features of A. R. Desai's Marxist sociology'' and notes Section A questions are static.
\end{PYQLink}

\begin{QuickRevision}
\begin{itemize}
\item Dominant caste = numbers + wealth + political power (three capitals).
\item Ritual hierarchy (purity) vs secular hierarchy (land/wealth/power); status is contextual.
\item Yadavs Sanskritized themselves to claim Kshatriya status.
\end{itemize}
\end{QuickRevision}

\section{Sanskritization Revisited (Conceptual Issues)}

The \textbf{original definition} of Sanskritization (later revised by Srinivas himself) was: ``a process in which low-caste Hindus orient themselves to the \textbf{Brahminical} lifestyle'' (initially called \textbf{Brahmanization}). Later he realised people follow different reference groups --- Kshatriya, Vaishya, even Shudra lifestyles --- so he revised it to ``imitation of the \textbf{dvija} lifestyle''.

\begin{ExampleBox}
In Gujarat the reference group is the \textbf{Patel} (merchant); in Haryana, the \textbf{Jat}; in Andhra, the \textbf{Reddy}; in U.P., the \textbf{Yadav}. All are \textbf{local jatis} --- none appears in the Varna hierarchy. Slogans: ``baniya saf, brahmin maaf'', ``brahmin baniya thakur chor, baaki sab DS4'' (DS4 = Dalit Shoshit Samaj, Kanshi Ram, 1981, U.P.).
\end{ExampleBox}

\begin{KeyConcept}
Sanskritization can be seen as \textbf{anticipatory socialization} (Merton). It has existed since ancient times --- e.g.\ the \textbf{Satavahana dynasty} (also called Andhra) stood out by claiming \textbf{Brahminical status} (a warrior class claiming Brahmin status).
\end{KeyConcept}

\begin{CriticalPoint}
\textbf{Desanskritization} (a term coined by D. N. Majumdar) is not academically well-accepted. \textbf{Rajputization} was coined by Surjit Sinha (among the Bhumij tribe of West Bengal).
\end{CriticalPoint}

\begin{QuickRevision}
\begin{itemize}
\item Original definition = Brahminical lifestyle (Brahmanization); revised to dvija lifestyle.
\item Reference groups vary (Patel, Jat, Reddy, Yadav --- all local jatis).
\item Sanskritization = anticipatory socialization; Satavahana claimed Brahminical status.
\end{itemize}
\end{QuickRevision}

\section{The Critique of Structural Functionalism}

\begin{itemize}
\item \textbf{Harmonious, rosy image}: SF only sees integration, not conflict, exploitation or untouchability. It negates material conditions (poverty, gender inequality). Srinivas saw the Oka (patrilineal) but not that a woman has no share in it.
\item \textbf{Brahminical perspective}: Srinivas (a Brahmin) reads through a Brahminical lens; ideology is justified in the name of sociology; his concepts assume everyone wants to claim upper-caste status.
\item \textbf{Sweeping generalizations}: from two field studies (Coorg, Rampura) he generalised to ``all India'' (inductive reasoning, not falsifiable). In the field, wealth/numbers/power are not neatly concentrated in one caste.
\item \textbf{Field studies reflect book-view biases}: despite knowing everyone claims superiority, he still believed all over India Brahmins are at the top and untouchables at the bottom (his Brahmin social location).
\end{itemize}

\begin{QuickRevision}
\begin{itemize}
\item SF critique: ignores conflict/material conditions; Brahminical bias; sweeping generalizations; book-view biases.
\item Provides a harmonious, not objective, picture.
\end{itemize}
\end{QuickRevision}
